% !TeX encoding = UTF-8
\documentclass[variant=guiaresuelta,cover=institutional,mode=screen]{ua-engineering-v6-article}
% !TeX encoding = UTF-8
\UASetUniversity{Universidad Austral}
\UASetFaculty{Facultad de Ingeniería}
\UASetDepartment{Departamento de Computación}
\UASetCampus{Pilar}
\UASetCareer{Ingeniería Informática}
\UASetCommission{Comisión A}
\UASetProfessor{Equipo docente}
\UASetSemester{Segundo cuatrimestre 2026}
\UASetCourse{Sistemas Distribuidos}
\UASetDate{2026}

\UASetClassNumber{09}
\UASetClassTitle{Arquitecturas distribuidas reales}
\title{Clase 09 --- Guía resuelta}
\author{\UAProfessor}
\date{\UADate}
\begin{document}
\maketitle
\section{Cómo leer las soluciones}
Son modelos de argumentación arquitectónica. Otras fronteras son válidas si compran autonomía real y declaran sus costos.
\section{Ejercicios y soluciones}
\begin{uaexercise}
\textbf{1. Monolito modular.} Explique por qué una unidad de despliegue no implica ausencia de modularidad. Proponga reglas internas para evitar una bola de barro.
\end{uaexercise}
\begin{uaanswer}
Un monolito puede imponer módulos con APIs internas, dependencias acíclicas, ownership, pruebas de arquitectura, esquemas o namespaces separados y prohibición de acceso lateral a internals. La bola de barro aparece cuando cualquier módulo modifica cualquier dato y las dependencias dejan de ser explícitas, no por compartir proceso.
\end{uaanswer}
\begin{uaexercise}
\textbf{2. Prima de microservicios.} Enumere los costos distribuidos que aparecen al extraer un módulo y diseñe un criterio para decidir si el beneficio compensa.
\end{uaexercise}
\begin{uaanswer}
Aparecen red, timeouts, retries, observabilidad, contratos, despliegue, secretos, datos distribuidos, sagas y on-call. Separar solo si existe presión real: escalado/ciclo de cambio/aislamiento/ownership distinto, y si una métrica demostrará mejora. Si el beneficio es solo 'arquitectura moderna', no compensa.
\end{uaanswer}
\begin{uaexercise}
\textbf{3. Señales de separación.} Dado un sistema con catálogo, pagos y recomendaciones, decida qué separar primero y qué mantener junto. Justifique por dominio, carga y equipo.
\end{uaexercise}
\begin{uaanswer}
Pagos puede separarse si requiere seguridad, compliance, equipo y carga distintos. Recomendaciones puede separarse por escalado y consistencia eventual. Catálogo quizá permanezca con núcleo si el dominio cambia. La decisión depende de fronteras maduras y capacidad operativa; no hay orden universal.
\end{uaanswer}
\begin{uaexercise}
\textbf{4. Monolito distribuido.} Construya una arquitectura con cinco servicios que conserve despliegue coordinado, base compartida y cadenas RPC. Explique por qué es peor que un monolito modular.
\end{uaexercise}
\begin{uaanswer}
Cinco servicios comparten DB, deben desplegar misma versión y se llaman en cadena. Se agregan latencia, fallas parciales y tracing, pero ningún equipo puede cambiar de forma autónoma. Un monolito modular sería más simple y ofrecería la misma coordinación interna sin red.
\end{uaanswer}
\begin{uaexercise}
\textbf{5. Bounded contexts.} Defina tres significados distintos de 'curso' o 'pedido' en contextos de negocio y diseñe traducciones entre ellos.
\end{uaexercise}
\begin{uaanswer}
Curso en catálogo: producto y descripción; en aprendizaje: módulos y progreso; en facturación: SKU, precio e impuesto. Los contratos traducen IDs y hechos relevantes; no se comparte un objeto universal. Esto permite que cada contexto evolucione sin imponer semántica ajena.
\end{uaanswer}
\begin{uaexercise}
\textbf{6. Ownership de datos.} Explique qué significa database per service y cómo resolvería una consulta global sin permitir acceso directo a tablas ajenas.
\end{uaexercise}
\begin{uaanswer}
El servicio es el único writer de su modelo. Una consulta global se resuelve con composición de APIs, materialized view alimentada por eventos, data warehouse o servicio de búsqueda. Elegir según frescura y latencia; acceso directo a tablas rompe ownership y contratos.
\end{uaanswer}
\begin{uaexercise}
\textbf{7. Conway.} Analice una organización donde dos equipos deben aprobar cada cambio del otro. ¿Qué revela sobre la frontera de servicios?
\end{uaexercise}
\begin{uaanswer}
La frontera técnica no coincide con autonomía organizacional. Si cada cambio requiere coordinación cruzada, los servicios forman una unidad de cambio. Puede convenir agrupar ownership, redefinir contratos o incluso mantener juntos los componentes hasta madurar la organización.
\end{uaanswer}
\begin{uaexercise}
\textbf{8. RPC vs eventos.} Para validación de pago, notificación, actualización de ranking y consulta de perfil, elija RPC o evento y justifique.
\end{uaexercise}
\begin{uaanswer}
Validación de pago suele requerir RPC sincrónico en checkout; notificación es evento; ranking es evento/proyección eventual; perfil puede ser RPC/cache según frescura. La elección se basa en necesidad de respuesta inmediata, tolerancia al atraso y ownership.
\end{uaanswer}
\begin{uaexercise}
\textbf{9. Cadena sincrónica.} Diseñe una cadena API->A->B->C y calcule disponibilidad compuesta si cada servicio tiene 99.9%. Analice además latencia y timeouts.
\end{uaexercise}
\begin{uaanswer}
Disponibilidad aproximada de la cadena es $0.999^4 \approx 99.6006\%$, sin contar red. Latencias se suman y colas multiplican p99. Timeouts/retries pueden amplificar carga. La cadena muestra que alta disponibilidad individual no garantiza flujo altamente disponible.
\end{uaanswer}
\begin{uaexercise}
\textbf{10. Comando, evento y consulta.} Clasifique ReserveSeat, SeatReserved y GetSeatAvailability. Construya un ejemplo de evento mal modelado como comando oculto.
\end{uaexercise}
\begin{uaanswer}
ReserveSeat es comando; SeatReserved evento; GetSeatAvailability consulta. Un 'PaymentRequested' publicado esperando que un único servicio obligatoriamente actúe y responda puede ser un comando disfrazado, con ownership y manejo de error ambiguos.
\end{uaanswer}
\begin{uaexercise}
\textbf{11. Dual write.} Un servicio actualiza Enrollment y publica StudentEnrolled. Dibuje las ventanas de falla para ambos órdenes.
\end{uaexercise}
\begin{uaanswer}
DB primero deja inscripción sin evento si crash; evento primero deja evento de inscripción que abortó. Si ambos ocurren pero ACK se pierde, retry puede duplicar. No existe orden que haga atómicas dos escrituras independientes.
\end{uaanswer}
\begin{uaexercise}
\textbf{12. Transactional outbox.} Diseñe schema mínimo, relay, deduplicación y política de retención para una outbox de inscripciones.
\end{uaexercise}
\begin{uaanswer}
Campos: \texttt{event\_id}, \texttt{aggregate\_id}, type, version, payload, \texttt{created\_at}, sequence, status. La transacción local inserta negocio+outbox. Relay publica repetidamente y marca; consumidores deduplican \texttt{event\_id}. Retención hasta confirmación y ventana de auditoría; monitorear edad y backlog.
\end{uaanswer}
\begin{uaexercise}
\textbf{13. Evolución de eventos.} Proponga una evolución compatible de StudentEnrolled v1 a v2 y una evolución que rompería consumidores.
\end{uaexercise}
\begin{uaanswer}
Compatible: agregar optional \texttt{scholarship\_code} con default ausente. Rompe: renombrar \texttt{student\_id} sin compatibilidad o cambiar amount de centavos a unidades manteniendo nombre. Debe existir versionado, schema validation y período de transición.
\end{uaanswer}
\begin{uaexercise}
\textbf{14. Saga de inscripción.} Diseñe pasos, timeouts, reintentos y compensaciones para inscripción, cupo, pago, acceso y comprobante.
\end{uaexercise}
\begin{uaanswer}
Crear inscripción pending; reservar cupo; autorizar pago; activar acceso; emitir comprobante. Retries idempotentes por saga/step ID. Si pago falla, liberar cupo y cancelar. Si acceso falla tras cobro, reintentar o refund según timeout. Email no debe revertir negocio.
\end{uaanswer}
\begin{uaexercise}
\textbf{15. Compensación imperfecta.} Dé tres efectos que no pueden revertirse perfectamente y proponga una compensación de negocio razonable.
\end{uaexercise}
\begin{uaanswer}
Email enviado: compensar con aclaración; envío físico: iniciar devolución; reserva aérea consumida: reembolso/credito, no borrar historia. La compensación es un nuevo efecto y debe modelar costo y observabilidad.
\end{uaanswer}
\begin{uaexercise}
\textbf{16. Coreografía.} Diseñe una saga coreografiada con tres participantes y explique cómo evitar ciclos y pérdida de visibilidad.
\end{uaexercise}
\begin{uaanswer}
Tres servicios reaccionan a eventos lineales y publican resultado. Evitar ciclos con contratos de workflow, ownership claro y pocos participantes; usar correlation/saga ID y una vista de seguimiento. Si crece, considerar orquestación.
\end{uaanswer}
\begin{uaexercise}
\textbf{17. Orquestación.} Diseñe el mismo flujo con un orquestador durable. Defina estado, recuperación y riesgo de centralización.
\end{uaexercise}
\begin{uaanswer}
El orquestador persiste estado por saga, paso, intentos, deadline y compensaciones. Al reiniciar continúa. Facilita timeout y visibilidad, pero se vuelve componente crítico y puede concentrar lógica de dominio; debe replicarse y versionarse.
\end{uaanswer}
\begin{uaexercise}
\textbf{18. Strangler Fig.} Proponga una migración incremental desde un monolito de e-commerce. Defina primer servicio, routing, rollback y métricas.
\end{uaexercise}
\begin{uaanswer}
Extraer primero una capacidad con frontera y presión real, por ejemplo recomendaciones. Gateway enruta porcentaje de tráfico. Comparar latencia/costo/error; rollback vuelve al monolito. Luego retirar código antiguo. No iniciar por el núcleo más acoplado.
\end{uaanswer}
\begin{uaexercise}
\textbf{19. Capacidad operativa.} Construya una checklist de capacidades mínimas antes de operar 30 microservicios en producción.
\end{uaexercise}
\begin{uaanswer}
CI/CD, rollback, infraestructura como código, service discovery, secretos, métricas/logs/traces, SLOs, alertas, on-call, catálogo/ownership, contratos, seguridad, resiliencia, costos, runbooks y pruebas de falla. Sin esto, 30 servicios multiplican toil.
\end{uaanswer}
\begin{uaexercise}
\textbf{20. Integración D=(P,F,G,S,T).} Aplique el marco a una plataforma de cursos y defienda qué permanece monolítico, qué se extrae y qué trade-offs se aceptan.
\end{uaexercise}
\begin{uaanswer}
P: inscribir y enseñar globalmente. F: fallas parciales, dual writes, picos, cambios de dominio. G: pago/cupo correctos, video disponible, recomendaciones frescas. S: monolito modular para aprendizaje cambiante, pagos/video separados, eventos para proyecciones, outbox y saga. T: más operación y consistencia eventual.
\end{uaanswer}
\section{Criterio docente de corrección}
La respuesta excelente identifica una presión concreta para separar, conserva invariantes locales y diseña operación. Es insuficiente responder con slogans de microservicios o eventos.
\end{document}
